\section{\textit{moeten} (deber, obligación)}
\textbf{Significado.} Expresa obligación interna o externa; se siente más fuerte que un simple consejo.

\textbf{Formas clave}
\begin{itemize}[leftmargin=*]
  \item \textit{ik moet}, \textit{jij moet}, \textit{wij/jullie/zij moeten}.
  \item Negar obligación inexistente suele usar \textit{hoeven \ldots{} niet} (vista previa): \textit{Ik hoef morgen niet te werken}.
\end{itemize}

\textbf{Ejemplos}
\begin{itemize}[leftmargin=*]
  \item \textit{Je moet je ticket tonen.} (Debes mostrar tu billete.)
  \item \textit{Wij moeten om zes uur klaar zijn.} (Tenemos que estar listos a las seis.)
  \item \textit{Ik hoef niet naar kantoor.} (No tengo que ir a la oficina.)
\end{itemize}

\textbf{Mini práctica}
\begin{itemize}[leftmargin=*]
  \item Completa: \textit{Jullie \_\_\_ (moeten) stil zijn}.
  \item Reformula en negativo con \textit{hoeven}: \textit{Ik moet koken}. Escribe la versión correcta.
\end{itemize}
